% =====================================================================================
% DATOS DEL ALUMNO Y DEL PROYECTO  -->  ESTE ES EL ÚNICO ARCHIVO DE CONFIGURACIÓN QUE DEBES EDITAR
% (además de escribir tus secciones en la carpeta secciones/)
% Cambia solo lo que está entre llaves { } a la derecha de cada \newcommand.
% El sistema de revisión lee estos datos (título, alumno, tipo de proyecto): no cambies los nombres
% de los comandos (\NombreAlumno, \NombreProyecto, \TipoProyecto, ...).
% =====================================================================================

% ---------- Tipo de proyecto (OBLIGATORIO) ------------------------------------------
% Escribe exactamente  sistema   o   investigacion
%   sistema        -> página web, sistema web, API, sistema de escritorio, app móvil
%   investigacion  -> algoritmo, modelo de aprendizaje de máquina u otro trabajo de investigación
% Según el tipo, la Revisión 2 y la 3 piden secciones distintas (ver GUIA_ALUMNO.md).
\newcommand{\TipoProyecto}{sistema}

% ---------- Estilo de citas y referencias -------------------------------------------
% ieee (por defecto) o apa. Pregunta a tu director cuál debes usar.
\newcommand{\EstiloCitas}{ieee}

% ---------- Alumno -------------------------------------------------------------------
\newcommand{\NombreAlumno}{Ezequiel Torres Desilos}
\newcommand{\Matricula}{2330231}
% Alumnas: pongan  la  y  a ;  alumnos: pongan  el  y  o
\newcommand{\elolaNombreAlumno}{el}
\newcommand{\OA}{o}

% ---------- Programa y proyecto ------------------------------------------------------
\newcommand{\ncarrera}{Ingeniería en Tecnologías de la Información}
\newcommand{\NombreProyecto}{Copiloto Agéntico RAG para Documentación Técnica y Procedimientos de Planta}
% Título para el encabezado de cada página. Usa \\ para partirlo en líneas (máximo 3 líneas).
\newcommand{\NombreProyectoheader}{Copiloto Agéntico RAG para Documentación \\Técnica y Procedimientos de Planta}

% ---------- Organismo receptor (empresa o institución donde realizas la estadía) -----
\newcommand{\organismoreceptor}{Centro de Innovación Industrial en Inteligencia Artificial (CII.IA)}
% Si el organismo es femenino (la universidad, la empresa) pon  la ; si es masculino (el instituto) pon  el
\newcommand{\elolaNombreEmpresa}{la}

% ---------- Asesores y evaluador -----------------------------------------------------
\newcommand{\nasesorinstitucional}{DR. SAID POLANCO MARTAGÓN}
\newcommand{\nasesorempresaria}{Dra. Ludivina Facundo Flores}
\newcommand{\nevalador}{Pendiente de asignación}

% ---------- Fechas -------------------------------------------------------------------
\newcommand{\ncuatrimestre}{septiembre-diciembre 2026}
\newcommand{\fechaPortada}{Diciembre de 2026}
\newcommand{\fechacarta}{0 de Mes de 2026}
\newcommand{\FechaExposicion}{0 de Mes de 2026}
\newcommand{\HoraExposionFormatoVenticuatroHoras}{0:00}
